\chapter{Electromagnetism}
\section{Debye-Huckel Theory}

Debye-Huckel theory is a description of the electromagnetic interactions between ions in a solution of ionic compounds.

\subsection{Poisson-Boltzmann}

The potential in a field of ions must be solved in terms of electrostatics in a media. This is described by Poisson's equation
\begin{equation}
    \nabla^2 \Psi(\mathbf{r}) = -4 \pi \frac{\rho(\mathbf{r})}{\varepsilon}
\end{equation}
In a sea of ions, we must carefully define the charge density as it pertains to the system. In a general solution of ions, there will be $j$ species of ions.
\begin{align*}
    \nabla^2 \Psi(\mathbf{r}) &=- 4 \pi \frac{\rho_j(\mathbf{r})}{\varepsilon}
\end{align*}
The charge distribution of a single ion concentration is essentially the density of the ionic positions. It would be tedious and unfeasible to count every ion's position exactly as we are assuming too many things. Rather, one can use a statistical ensemble of ions via a cannonical ensemble, which is a good approximation here.

\begin{equation}
    \rho_i(\mathbf{r}) = C_i e_0 Z_i \exp{\left[-\beta Z_i e_0\Psi(r)\right]}
\end{equation}

Here the energy of the system is defined as the potential energy scaled by the charge defined as $U = qV$. The full equation is known as the Poisson-Boltzmann equation, where $\Psi(\mathbf{r})$ is the potential for all space.
\begin{equation}
    \nabla^2 \Psi(\mathbf{r}) = -\frac{1}{\varepsilon_0 \varepsilon_r}\sum_{i =1}^\infty C_i e_0 Z_i \exp{\left[-\beta Z_i e_0\Psi(r)\right]}
\end{equation}

There is a nice approximation we can make given that the potential of the solution is small. We linearize the exponential to first order here.

\begin{align*}
    \nabla^2 \Psi(\mathbf{r}) &\approx -\frac{1}{\varepsilon_0 \varepsilon_r}\sum_{i =1}^\infty C_i e_0 Z_i \left[1-\beta Z_i e_0\Psi(r)\right] \\
    &= -\frac{1}{\varepsilon_0 \varepsilon_r}\sum_{i =1}^\infty \left[C_i e_0 Z_i-\beta C_i Z_i^2 e_0^2\Psi(r)\right]\\
    &= -\sum_{i =1}^\infty \left[\frac{C_i e_0 Z_i}{\varepsilon_0 \varepsilon_r}-\kappa^2 \Psi(r)\right]
\end{align*}
Where $\kappa_i^2$ is related to the debye-length, a unit of the electrostatic screening caused by the ionic interactions. The Debye length is defined as $\lambda_D = \kappa^{-1}$.

\begin{equation}
 \kappa_i^2 = \frac{\beta C_i Z_i^2 e_0^2}{\varepsilon_0 \varepsilon_r}
\end{equation}

\subsection{Monovalent Ionic Mixture}
Suppose there is a single type of ion in the mixture. The sum collapses into a single term and the constant term is dropped.
\begin{equation}
    \nabla^2 \Psi(\mathbf{r}) = \kappa^2 \Psi(\mathbf{r})
\end{equation}
This equation has an obvious solution which are the real exponentials;
\begin{equation}
    \Psi(\mathbf{r}) = A e^{-\kappa\mathbf{r}} + B e^{\kappa\mathbf{r}}
\end{equation}
